\section{十、从门级网表到可制造芯片}

综合完成时，设计仍然只是一张带时序约束的门级网表：其中已经出现具体的与非门、触发器、时钟门控单元和存储宏，却还没有决定这些单元位于硅片的什么位置，也没有形成可供制造的几何图形。把网表变成芯片，需要后端实现、签核、制造、封装和首次上电五个连续阶段。每个阶段都产生一种新的事实；后一个阶段不能用自己的通过结果替代前一个阶段的证明。

\begin{pdfdiagram}[x=1cm,y=1cm]
  \node[hw state, minimum width=2.5cm] (rtl) at (1.4,2.4) {RTL 与约束};
  \node[hw compute, minimum width=2.5cm] (syn) at (4.6,2.4) {综合与映射};
  \node[hw state, minimum width=2.5cm] (net) at (7.8,2.4) {门级网表};
  \node[hw compute, minimum width=2.5cm] (pnr) at (10.9,2.4) {布局布线};
  \draw[hw bus] (rtl) -- node[hw label,above]{逻辑/时序约束} (syn);
  \draw[hw bus] (syn) -- node[hw label,above]{库单元实例} (net);
  \draw[hw bus] (net) -- node[hw label,above]{位置与连线} (pnr);

  \node[hw control, minimum width=2.5cm] (sign) at (10.9,0.2) {签核与版图};
  \node[hw block, minimum width=2.5cm] (fab) at (7.8,0.2) {掩膜与晶圆制造};
  \node[hw block, minimum width=2.5cm] (pkg) at (4.6,0.2) {封装与量产测试};
  \node[hw state, minimum width=2.5cm] (bring) at (1.4,0.2) {上电与硅后验证};
  \draw[hw bus] (pnr) -- (10.9,1.3) -- (sign);
  \draw[hw bus] (sign) -- node[hw label,above]{GDS/掩膜数据} (fab);
  \draw[hw bus] (fab) -- node[hw label,above]{合格裸片} (pkg);
  \draw[hw bus] (pkg) -- node[hw label,above]{可供电器件} (bring);
\end{pdfdiagram}
\diagramnote{RTL 到可运行芯片的两行实现链。上行把行为约束落实为单元位置和物理连线；下行把通过签核的版图变成晶圆、封装器件和可上电样片。箭头上的产物是阶段间的交接契约，不能把“仿真通过”直接解释为“芯片可制造”。}

\topic{标准单元库把逻辑功能连接到器件与版图}

\term{标准单元}{standard cell}是预先完成晶体管级设计、版图和时序表征的基本块。一个二输入与非门通常同时提供多种驱动强度和阈值电压版本：高驱动版本能推动更大负载，但面积与动态功耗更高；低阈值版本切换更快，但漏电更大。综合器先选择逻辑功能，后端工具再根据扇出、线长、时序和功耗不断调整具体版本。存储器通常不由标准单元拼接，而以 SRAM、ROM 或寄存器文件宏的形式进入平面规划，因为宏的端口、宽高比、电源引脚和内部时序已经固定。

\term{平面规划}{floorplan}先确定芯片边界、I/O 区域、电源网格和大宏位置。宏若挡住主要数据通路，后续布线会产生绕行；宏若离使用者过远，访问延迟和动态功耗会上升；电源网格过弱则会在同时切换时产生 IR drop。平面规划不是美观问题，而是把面积、连线、电源和热密度提前放进同一张约束图。

\topic{布局、时钟树与布线把网表变成有延迟的几何结构}

\term{布局}{placement}为每个标准单元选择合法位置。工具在满足行高、阻塞区和拥塞约束的同时，尽量缩短关键网线并控制局部功耗密度。布局之后进行\term{时钟树综合}{Clock Tree Synthesis, CTS}：从时钟入口插入分级缓冲和必要的门控单元，使各寄存器看到的时钟沿满足偏斜、转换时间和负载限制。时钟树完成后，保持时间可能比布局阶段更紧，因为真实时钟到达差已经出现。

\term{布线}{routing}为电源、时钟和普通信号分配金属层、线宽、间距与过孔。局部低层金属连接相邻单元，高层宽线承担长距离、时钟和电源。布线产生的电阻与电容会反过来改变门延迟，因此工具必须抽取寄生参数，再用更新后的延迟重新执行静态时序分析。若一条路径失败，修复可能是换单元、插缓冲、移动逻辑、调整时钟或增加流水级；任何修复都要重新检查其他模式和工艺角。

\topic{签核回答“可制造且在规定边界内工作”}

\begin{longtable}{L{0.18\linewidth}L{0.34\linewidth}L{0.38\linewidth}}
\toprule
签核项目 & 检查的物理事实 & 失败意味着什么 \\
\midrule
STA & 所有时钟模式、PVT 角与例外路径 & 数据可能在采样边界前未稳定，或旧值保持时间不足 \\\tablerowrule
DRC & 线宽、间距、包围、密度等制造规则 & 几何图形超出工艺可稳定制造的窗口 \\\tablerowrule
LVS & 版图抽取网表与设计网表是否等价 & 版图中存在漏连、短接或器件参数不一致 \\\tablerowrule
功耗与 IR drop & 电源网格上的电流、压降和热点 & 单元在负载瞬态下可能得到不足电压 \\\tablerowrule
电迁移 & 金属与过孔的长期电流密度 & 连线可能在寿命期内因原子迁移失效 \\\tablerowrule
信号完整性 & 串扰、转换时间与噪声窗口 & 相邻活动可能改变关键网的延迟或逻辑判决 \\
\bottomrule
\end{longtable}

签核使用的是特定工艺设计套件、单元库、寄生模型和工作条件。通过签核表示“在这些模型和角落内没有发现违反约束”，不表示所有样片必然完全相同。制造波动、封装寄生、板级供电和未建模使用条件仍要由量产测试与硅后验证继续约束。

\topic{从掩膜到封装，逻辑状态变成可供电的器件}

签核版图经数据准备生成各层掩膜。晶圆制造反复执行薄膜沉积、光刻、刻蚀、离子注入和平坦化，在硅上形成晶体管与多层互连。制造完成后先做晶圆级探针测试，把明显失效的裸片排除；合格裸片切割、贴装并通过键合线、倒装焊凸点或硅中介层连接到封装基板。封装同时承担电源分配、信号逃逸、机械保护和热传导，封装寄生电感与热阻会进入真实工作边界。

首次上电应从最小可信条件开始：确认电源轨顺序与静态电流，确认复位与参考时钟，再尝试 JTAG 身份读取、片内 SRAM 测试、Boot ROM 取指和外部接口训练。若一上电就运行复杂软件，任何异常都可能同时来自电源、时钟、制造缺陷、封装连接、DFT 配置或逻辑设计，证据难以分层。

\section{十一、FPGA：位流怎样把通用阵列变成专用电路}

FPGA 不是“可以运行 Verilog 的处理器”。它是一张预制的可配置硬件网格，配置位决定查找表实现哪条真值关系、交换矩阵接通哪些导线、触发器是否启用、I/O 使用何种电平标准。RTL 综合与布局布线完成后生成的\term{位流}{bitstream}写入配置存储器，配置完成时，整张通用网格同时成为目标电路。

\begin{pdfdiagram}[x=1cm,y=1cm]
  \node[hw control, minimum width=2.4cm] (cfg) at (1.4,3.2) {配置端口\\位流装载};
  \node[hw state, minimum width=2.6cm] (bits) at (4.5,3.2) {配置 SRAM\\连接与功能位};
  \draw[hw bus] (cfg) -- (bits);

  \node[hw compute, minimum width=2.5cm, minimum height=1.2cm] (clb) at (2.0,1.0) {逻辑块\\LUT + FF + carry};
  \node[hw compute, minimum width=2.5cm, minimum height=1.2cm] (dsp) at (5.1,1.0) {DSP 阵列\\乘加与累加};
  \node[hw memory, minimum width=2.5cm, minimum height=1.2cm] (bram) at (8.2,1.0) {BRAM/URAM\\片上存储};
  \node[hw block, minimum width=2.5cm, minimum height=1.2cm] (io) at (11.3,1.0) {I/O + SerDes\\PLL/MMCM};
  \draw[hw bus] (clb) -- node[hw label,above]{可编程互连} (dsp);
  \draw[hw bus] (dsp) -- node[hw label,above]{可编程互连} (bram);
  \draw[hw bus] (bram) -- node[hw label,above]{可编程互连} (io);
  \draw[hw return] (bits.south) -- (4.5,2.25) -- (6.6,2.25);
  \node[hw label,anchor=west] at (6.7,2.25) {控制 LUT 内容、开关矩阵、I/O 与时钟资源};
\end{pdfdiagram}
\diagramnote{FPGA 的两层结构。上层配置存储器保存位流状态，下层是被这些状态重新连接的逻辑、DSP、存储、I/O 和时钟资源。位流不是程序指令流；它在配置边界上一次性决定后续每个时钟周期中并行存在的电路。}

\topic{LUT、触发器和进位链构成基本逻辑块}

$k$ 输入 LUT 本质上是一块 $2^k$ bit 的小存储器：输入作为地址，存储位作为输出，因此可以表达任意 $k$ 输入布尔函数。LUT 后接触发器，允许同一逻辑块输出组合结果或寄存结果。专用 carry chain 把相邻逻辑块的进位直接串联，避免宽加法器经过普通交换矩阵；这就是 FPGA 上加法器通常比任意同深度 LUT 网络更快的原因。

普通可编程互连需要配置开关和长线段，延迟常常高于 LUT 本身。FPGA 时序优化因而不只减少逻辑级数，还要减少远距离跨区、控制高扇出并使用专用时钟网络。BRAM、DSP、SerDes 和 PLL 是硬化资源：把常用结构直接做成固定电路，单位面积、速度和功耗显著优于用 LUT 拼接，但端口形状与数量受到器件限制。

\topic{配置启动是一条独立于用户逻辑的状态机}

上电后，FPGA 先保持用户 I/O 在安全状态，清除或初始化配置存储器，再从 SPI flash、并行 flash、JTAG 或主机接口接收位流。配置逻辑逐帧校验并写入内部配置 SRAM，全部完成后才释放全局启动信号，让用户时钟、触发器和 I/O 进入运行态。若位流校验失败、供电不稳定或配置时钟停止，用户逻辑不会获得一个“部分可用”的正常状态；平台应保持输出安全并报告配置错误。

位流还包含器件型号、资源位置和布线开关选择，通常不能跨器件直接复用。RTL 是相对可移植的功能描述，约束与位流则属于具体板卡和具体器件。调试时必须区分三类失败：RTL 行为错误、实现时序失败、配置与板级边界失败。

\section{十二、DFT 与量产测试：怎样观察芯片内部状态}

芯片正常工作时，数百万个触发器只通过少量外部引脚间接影响结果。制造测试若仍从功能接口施加输入，许多内部节点需要运行很长序列才能被激活和观察。\term{可测性设计}{Design for Test, DFT}增加专用测试路径，把“难以控制、难以观察”的内部状态转换成可移入和可移出的结构。

\begin{pdfdiagram}[x=1cm,y=1cm]
  \node[hw block, minimum width=2.4cm] (si) at (1.4,2.6) {scan in};
  \node[hw state, minimum width=2.2cm] (ff0) at (4.0,2.6) {扫描触发器 0};
  \node[hw state, minimum width=2.2cm] (ff1) at (6.8,2.6) {扫描触发器 1};
  \node[hw state, minimum width=2.2cm] (ff2) at (9.6,2.6) {扫描触发器 2};
  \node[hw block, minimum width=2.4cm] (so) at (12.2,2.6) {scan out};
  \draw[hw bus] (si) -- (ff0);
  \draw[hw bus] (ff0) -- (ff1);
  \draw[hw bus] (ff1) -- (ff2);
  \draw[hw bus] (ff2) -- (so);

  \node[hw compute, minimum width=2.6cm] (comb0) at (5.4,0.4) {功能组合逻辑 A};
  \node[hw compute, minimum width=2.6cm] (comb1) at (8.2,0.4) {功能组合逻辑 B};
  \draw[hw return] (ff0.south) -- (4.0,1.45) -- (comb0.west);
  \draw[hw return] (comb0.east) -- (6.8,1.45) -- (ff1.south);
  \draw[hw return] (ff1.south east) -- (7.1,1.2) -- (comb1.west);
  \draw[hw return] (comb1.east) -- (9.6,1.2) -- (ff2.south);
  \node[hw label] at (6.8,3.55) {scan\_enable=1：移位装载/读出；scan\_enable=0：按功能路径采样};
\end{pdfdiagram}
\diagramnote{同一组三只扫描触发器的两种连接视图。测试模式把触发器串成可移位链，先装入指定内部状态，再切回功能模式捕获一拍组合逻辑响应，最后切回扫描模式读出。两种模式共享存储元件，但具有不同的数据路径和完成边界。}

\topic{scan 与 ATPG把内部组合逻辑变成可直接施测的对象}

扫描触发器在功能模式下与普通触发器相同；测试模式下，多只触发器首尾相连形成 scan chain。测试机先移入一组内部状态，切回功能时钟捕获一拍响应，再把结果移出。\term{自动测试向量生成}{Automatic Test Pattern Generation, ATPG}根据网表和故障模型寻找能激活并观察目标故障的向量。最常见的 stuck-at 模型假设节点固定为 0 或 1；transition 模型检查节点能否在规定时间内完成跳变。

大型芯片会并行使用多条扫描链并配合压缩器，缩短移位时间和外部引脚需求。扫描模式会改变时钟、功耗和时序条件，因此测试控制器必须规定移位频率、捕获脉冲、复位状态与模拟宏隔离，不能把测试模式视为普通功能模式的另一组输入。

\topic{存储器与高速接口需要专用自测}

SRAM 宏内部节点无法由普通 scan 直接覆盖，通常使用\term{存储器内建自测}{Memory Built-In Self-Test, MBIST}控制器执行 March 类序列：按正向和反向地址顺序进行写 0、读 0、写 1、读 1，暴露单元、耦合、译码和读写路径故障。发现坏行或坏列后，冗余分析逻辑选择 spare row/column，并把修复映射写入熔丝或一次性存储。

\term{逻辑内建自测}{Logic BIST, LBIST}常以 LFSR 产生伪随机激励，以多输入签名寄存器压缩响应。它适合上电自检和现场诊断，但压缩签名只能说明结果是否匹配，不能单独定位每个内部故障。SerDes、PLL、ADC 等模拟或混合信号宏还需要环回、眼图监测、校准码和专用模拟测试入口。

\topic{JTAG 边界扫描把芯片内部测试延伸到电路板}

JTAG TAP 由 TCK、TMS、TDI、TDO 和可选 TRST 构成状态机。边界扫描寄存器位于核心逻辑与封装引脚之间：正常模式透明传递信号，测试模式则可控制输出引脚并采样输入引脚。这样不运行 CPU、不依赖 DRAM，也能检查 BGA 封装下不可接触的焊点、相邻引脚短路和板级开路。JTAG 还可选择芯片 ID、旁路寄存器和片内调试模块，但量产产品必须用认证、生命周期状态或熔丝策略限制调试访问。

\topic{量产测试区分结构缺陷、速度分档与现场可靠性}

\begin{longtable}{L{0.18\linewidth}L{0.28\linewidth}L{0.42\linewidth}}
\toprule
阶段 & 主要对象 & 通过边界 \\
\midrule
晶圆探针测试 & 裸片电源、scan、MBIST、基础模拟宏 & 排除明显缺陷并记录可修复单元 \\\tablerowrule
封装后测试 & 封装连接、全速功能、温度与接口 & 确认封装寄生和装配未破坏功能 \\\tablerowrule
速度/功耗分档 & 最大频率、漏电、工作电压 & 把同一设计按可稳定工作范围分成产品等级 \\\tablerowrule
老化与抽检 & 高温高压、热循环和长期应力 & 估计早期失效率并验证寿命模型 \\\tablerowrule
系统上电自检 & 现场 SRAM、链路、时钟与传感器 & 确认当前机器满足继续启动的最低条件 \\
\bottomrule
\end{longtable}

量产测试通过不表示芯片在任意板卡和任意固件下都可运行；它证明的是器件在规定测试条件内满足结构、速度和功耗界限。系统启动仍要重新验证供电、时钟、存储训练和外部链路，因为这些边界直到芯片被装入具体机器后才形成。
